% AIRI Summer 2026 --- отчёт по проекту (ЧЕРНОВИК)
% Компилируется на Overleaf через pdfLaTeX.
\documentclass[10pt,a4paper]{article}

\usepackage[utf8]{inputenc}
\usepackage[T2A]{fontenc}
\usepackage[russian]{babel}
\usepackage[margin=2cm]{geometry}
\usepackage{booktabs}
\usepackage{array}
\usepackage{makecell}
\usepackage{graphicx}
\usepackage{xcolor}
\usepackage{hyperref}
\usepackage{caption}
\hypersetup{colorlinks=true,linkcolor=blue,urlcolor=blue}

\setlength{\tabcolsep}{4pt}
\renewcommand\theadfont{\normalsize}

\title{\vspace{-1.5em}LLM-Guided Performance Engineering for GPU Kernels}
\author{AIRI Summer 2026}
\date{\today}

\begin{document}
\maketitle
\vspace{-1.5em}

\section{Постановка задачи}

Проект исследует два аспекта построения LLM-агента, который разрабатывает и оптимизирует
GPU-ядра внутри фреймворка, плохо знакомого модели.

Первое направление --- освоение нового фреймворка прямо в процессе работы агента. Модели
неплохо справляются с Triton, но испытывают трудности с менее распространёнными
абстракциями, такими как CuTe. Чтобы понять, чем компенсировать этот пробел, нужно
окружение, подающее модели документацию, примеры кода, результаты компиляции, тесты и
данные профилирования; на нём уже можно сравнивать способы подачи информации и виды
обратной связи. Ожидаемый результат --- анализ инструментов, позволяющих LLM освоить новый
GPU-фреймворк \emph{без дообучения}.

Второе направление --- незнакомые числовые форматы на примере FP8: переход с FP16 требует
учёта специфики кодирования и специализированных инструкций, поэтому нужна отдельная
обвязка для генерации FP8-кода и устойчивый замкнутый цикл разработки. Оба направления мы
измеряем одними и теми же метриками: корректность, ускорение, число итераций агента и
стоимость в токенах.

\section{Инфраструктура эксперимента}

\paragraph{Модель.} Мы задеплоили Qwen3.6-35B-A3B в квантизации Q8\_0 на сервере с
2$\times$V100 с пайплайн-параллелизмом ($\mathrm{PP}=2$) и открыли доступ через
OpenAI-совместимый эндпоинт. В последствии, мы перешли на внешнего провайдера - DeepSeek v4 Flash, так как Qwen совсем не умел справляться с задачей. 

\paragraph{Оценщик.} Кандидатное ядро --- единственный файл, который агент имеет право
править. Всё остальное принадлежит оценщику: генерация входных данных, эталонная
реализация на Torch, проверка допусков, замер времени и вердикт \texttt{PASS}. Агент не
может подменить проверку и не видит эталонных таймингов. Ядра исполняются удалённо на
NVIDIA B300; для стабильности замеров делается несколько разминочных прогонов, затем
несколько основных с усреднением. Перед отправкой на GPU кандидат проходит локальную
статическую проверку \texttt{cute-check} (синтаксис и типы, без обращения к устройству)
--- она дёшево отсеивает заведомо нерабочий код.

\paragraph{Задачи.} Готовых бенчмарков по CuTe DSL и FP8 не существует, поэтому мы
адаптировали задачи из KernelBench: переиспользовали формат, но написали собственные
эталонные решения, привели входы и эпилоги к FP8-семантике и зафиксировали свои пороги
численного расхождения. Всего девять задач --- первого уровня (одиночная операция) и
второго (слитая последовательность: GEMM вместе с эпилогом). Эксперименты E1--E4
прогонялись на одной задаче второго уровня.

\paragraph{Справочные материалы.} Документацию по CuTe DSL мы разбили на четыре
кумулятивных уровня, не зависящих от конкретной задачи: \emph{только задача} (постановка и
стартовый файл), \emph{+ документация} (CuTe DSL API, алгебра layout'ов, архитектура
Blackwell, TMA и TMEM, численные свойства FP8/FP4), \emph{+ примеры} (нейтральные фрагменты
рабочего кода) и \emph{+ разбор ошибок} (типовые классы ошибок и то, как они выглядят в
диагностике). Каждый уровень включает предыдущий целиком, поэтому разница между соседними
строками таблиц --- ровно вклад добавленного материала.

\section{Как устроены эксперименты E1--E4}

Всего мы прогнали семь конфигураций. Первые четыре --- лестница уровней справки, они
отвечают на вопрос, \emph{какое знание вообще работает} (E1). Оставшиеся три меняют ровно
один параметр относительно лучшего уровня лестницы: способ доставки того же материала
(E2), наличие второго агента-критика между ходами (E3) и возврат сводки профилировщика
следующему ходу (E4). Конфигурация с разбором ошибок из E1 служит для них общей опорной
точкой, поэтому каждое сравнение однопараметрическое.

Каждая конфигурация --- один независимый запуск из шести \emph{ходов}. Ход --- итерация
работы агента: он пишет или правит кандидатное ядро, прогоняет его через статическую
проверку и через GPU, получает диагностику и переходит к следующему ходу. Эталонная
реализация перезамеряется внутри каждой конфигурации (0.3815--0.3835 мс), чтобы дрейф
загрузки машины не искажал сравнение ускорений.

\paragraph{Как читать таблицы.} «PASS $k/6$» --- в скольких из шести ходов получилось
корректное ядро. «Лучшее» --- время
самого быстрого корректного ядра за запуск, «ускорение» --- во сколько раз оно быстрее
эталона. Вход разбит на общий и попавший в кэш. Каждая конфигурация прогонялась один
раз.

\section{E1 --- какое знание важно}

Лестница из четырёх уровней: от голой постановки задачи (195 токенов справки) до полного
стека с разбором типовых ошибок (33\,146 токенов). Вопрос --- где проходит граница между
«модель не может» и «модель может».

\begin{table}[h]
\centering\small
\caption{E1, лестница уровней справки. Шесть ходов на уровень, время в мс.}
\label{tab:e1}
\resizebox{\textwidth}{!}{%
\begin{tabular}{lrcccrrrrrrr}
\toprule
Уровень справки & \thead{Справка,\\токенов} & PASS & \thead{Ходов до\\1-го PASS} & \thead{Лучшее,\\мс} & Ускор. & Вход & \thead{из них\\из кэша} & Выход & \thead{Токены\\рассужд.} & \thead{Токенов до\\1-го PASS} & \thead{Проверок\\на GPU} \\
\midrule
Только задача    & 195      & 0/6 & --- & ---    & ---   & 10{,}323{,}585 & 9{,}936{,}512  & 61{,}096 & 198{,}552 & ---           & 8  \\
+ документация   & 17{,}860 & 0/6 & --- & ---    & ---   & 14{,}519{,}559 & 13{,}687{,}552 & 92{,}626 & 470{,}913 & ---           & 15 \\
+ примеры        & 26{,}416 & \textbf{5/6} & 1 & 0.1181 & 3.23$\times$ & 8{,}762{,}557 & 8{,}308{,}864 & 72{,}828 & 231{,}642 & 2{,}805{,}178 & 25 \\
+ разбор ошибок  & 33{,}146 & \textbf{6/6} & 1 & \textbf{0.1160} & \textbf{3.29$\times$} & 7{,}434{,}221 & 6{,}977{,}792 & 60{,}662 & 179{,}156 & 1{,}604{,}327 & 14 \\
\bottomrule
\end{tabular}}
\end{table}

Граница проходит между документацией и примерами, и это не тот ответ, который подсказывает
объём текста. Почти восемнадцать тысяч токенов связного описания API, алгебры layout'ов и
архитектуры Blackwell не сдвигают результат ни на один ход: 0/6, ровно как на голой
постановке. При этом конфигурация с документацией --- \emph{самая дорогая} из семи:
14.5\,млн входных токенов и 470\,913 токенов рассуждений, вдвое с лишним больше, чем на
голой задаче. Модель добросовестно читает, думает дольше всех, пишет больше всех --- и не
решает ничего. Вредна она вполне определённым образом: ни одна из 21 проверки не прошла, и
восемь отказов --- это обращения к функциям, которых в библиотеке нет. Прочитав описание
API прозой, модель приобретает словарь, но не приобретает синтаксиса. Без справки поведение
другое: там кандидат до устройства просто не доходил. Без документации модель сдаётся, с
документацией --- галлюцинирует.

Перелом даёт добавление примеров: 5/6 и первое корректное ядро уже на первом ходу,
ускорение 3.23$\times$. Закрепляет его разбор ошибок --- 6/6, все 20 проверок пройдены, ни
одного отказа за запуск, лучшее ядро 0.1160\,мс (3.29$\times$). И лучший уровень оказался
\emph{самым дешёвым}: при максимальном объёме справки он потратил 7.43\,млн входных токенов
вместо 10.32\,млн и 179\,156 токенов рассуждений --- меньше, чем любая другая конфигурация;
до первого решения дошёл за 1.60\,млн токенов против 2.81\,млн у уровня с примерами. Вывод
E1: полезность справки не пропорциональна её объёму и даже не растёт вместе с ним. Работают
конкретные артефакты --- рабочие фрагменты кода и разбор типовых ошибок; описательные главы
про API не работают вовсе, но исправно тратят бюджет. Дальше уровень с разбором ошибок
используется как опорная конфигурация.

\section{E2 --- как доставлять справку}

E2 сравнивает два способа доставки \emph{одного и того же} материала. В опорной
конфигурации (\texttt{files}) файлы справки лежат в рабочей директории, и агент открывает
их сам, когда сочтёт нужным. В альтернативе (\texttt{prompt}) весь материал целиком
вставляется в промпт на каждом ходу. Содержимое идентично, различается только то, кто
решает, что читать.

\begin{table}[h]
\centering\small
\caption{E2, доставка справки: файлами в рабочей директории против вставки в промпт.}
\label{tab:e2}
\resizebox{\textwidth}{!}{%
\begin{tabular}{llcccrrrrrc}
\toprule
Конфигурация & Доставка & PASS & \thead{Ходов до\\1-го PASS} & \thead{Лучшее,\\мс} & Ускор. & Вход & \thead{\textbf{Вход на одно}\\\textbf{решение}} & Выход & \thead{Токенов до\\1-го PASS} & Таймауты \\
\midrule
Разбор ошибок (E1) & \texttt{files}  & \textbf{6/6} & 1 & 0.1160 & 3.29$\times$ & 7{,}434{,}221  & \textbf{1{,}239{,}037} & 60{,}662  & 1{,}604{,}327 & 0 \\
E2                 & \texttt{prompt} & 4/6 & 3 & \textbf{0.1080} & \textbf{3.54$\times$} & 12{,}986{,}542 & 3{,}246{,}636 & 102{,}068 & 5{,}648{,}113 & \textbf{3} \\
\bottomrule
\end{tabular}}
\end{table}

Результат расщепляется по метрикам. По пиковой производительности выигрывает вставка в
промпт: 0.1080\,мс и 3.54$\times$ --- лучшее ядро за все семь конфигураций. По всему
остальному она проигрывает: устойчивость падает с 6/6 до 4/6, первое корректное ядро
появляется только на третьем ходу, и это единственная конфигурация с таймаутами --- три
хода из шести исчерпали бюджет, не дойдя до проверки.
Постоянное присутствие всего материала не даёт агенту «забыть» редкие конструкции, отсюда
более высокий потолок оптимизации; но тот же материал вытесняет из внимания рабочий
контекст --- собственный код, диагностику компилятора, историю ходов. Косвенно это
подтверждает характер отказов: обращений к несуществующему API нет ни одного (документация
всегда перед глазами), зато есть пять ошибок компиляции --- модель знает, \emph{что}
писать, но теряет нить того, \emph{что она уже написала}.

Экономическая сторона однозначна: вход на одно полученное решение --- 1.24\,млн против
3.25\,млн токенов, доставка файлами дешевле в 2.6 раза; до первого решения разрыв ещё
больше, 1.60\,млн против 5.65\,млн. Практический вывод: чтение справки агентом по
требованию --- разумный режим по умолчанию, а вставку в промпт имеет смысл рассматривать
точечно, на финальной стадии оптимизации уже работающего ядра, где пиковая скорость важнее
устойчивости и цены.

\section{E3 --- обратная связь от агента-критика}

E3 добавляет между ходами вызов второго агента --- критика, работающего только на чтение.
Он получает текущее кандидатное ядро, диагностику с B300 и тот же уровень справки, что и
основной агент, и возвращает короткие подсказки. Сам он ничего не исполняет и не правит.
Его токены учитываются отдельно --- иначе сравнение стоимости было бы нечестным.

\begin{table}[h]
\centering\small
\caption{E3, агент-критик между ходами. Токены критика учтены отдельно.}
\label{tab:e3}
\resizebox{\textwidth}{!}{%
\begin{tabular}{lcccrrrrc}
\toprule
Конфигурация & PASS & \thead{Ходов до\\1-го PASS} & \thead{Лучшее,\\мс} & Ускор. & \thead{Вход\\осн. агента} & \thead{Выход\\осн. агента} & \thead{\textbf{Токены}\\\textbf{критика}} & Подсказок \\
\midrule
Разбор ошибок (E1) & \textbf{6/6} & 1 & \textbf{0.1160} & 3.29$\times$ & 7{,}434{,}221  & 60{,}662 & ---              & --- \\
E3                 & 4/6 & 1 & 0.1164 & 3.29$\times$ & 12{,}519{,}376 & 96{,}038 & \textbf{100{,}722} & 13 \\
\bottomrule
\end{tabular}}
\end{table}

Эффект нулевой по скорости и отрицательный по всему остальному. Лучшее ядро --- 0.1164
против 0.1160\,мс, то есть то же ускорение 3.29$\times$ с точностью до шума. Устойчивость
падает с 6/6 до 4/6. Вход основного агента вырастает на 68\,\%, выход --- на 58\,\%, и
сверху добавляются 100\,722 токена критика на 13 подсказок, около 7\,750 токенов на
подсказку. Мы платим примерно в полтора раза больше и получаем менее надёжный запуск.

Диагноз виден в характере отказов: из 11 неудачных проверок \textbf{десять --- это
обращения к несуществующему API}. Критик воспроизводит ровно ту патологию, которую мы
видели в E1: он читает документацию, но ничего не исполняет, поэтому его подсказки
правдоподобны по форме и называют функции, которых нет. Основной агент доверяет им больше,
чем собственной обратной связи от компилятора: подсказка приходит как внешний авторитет, а
сообщение об ошибке нужно ещё интерпретировать. Вывод: советчик без исполнения --- это
дополнительный источник галлюцинаций, а не фильтр от них. Если возвращаться к идее
критика, у него должно быть право самому запустить \texttt{cute-check} или пробную
компиляцию и высказываться только о проверенных гипотезах.

\section{E4 --- обратная связь от профилировщика}

E4 возвращает следующему ходу компактную сводку по времени исполнения участков ядра,
снятую на B300. Профилировать можно только работающее ядро, поэтому сводка появляется лишь
\emph{после} первого успеха: по построению она не может помочь пройти тест впервые, и
ожидаемый эффект здесь --- на ускорение среди уже проходящих ходов.

\begin{table}[h]
\centering\small
\caption{E4, возврат сводки профилировщика следующему ходу.}
\label{tab:e4}
\begin{tabular}{lccrrrr}
\toprule
Конфигурация & PASS & \thead{Лучшее,\\мс} & Ускор. & Вход & Выход & \thead{Проверок\\на GPU} \\
\midrule
Разбор ошибок (E1) & 6/6 & 0.1160 & 3.29$\times$ & 7{,}434{,}221  & 60{,}662 & 14 \\
E4                 & 5/6 & \textbf{0.1133} & \textbf{3.38$\times$} & 11{,}168{,}385 & 83{,}196 & 19 \\
\bottomrule
\end{tabular}
\end{table}

Так и вышло: доля успешных ходов практически не меняется (5/6 против 6/6), а лучшее ядро
улучшается с 0.1160 до 0.1133\,мс --- с 3.29$\times$ до 3.38$\times$, около 2.3\,\% по
времени. Качество кода при этом ровнее: за весь запуск всего два отказа, оба на компиляции,
ни одного обращения к несуществующему API и ни одного численного расхождения.

\begin{table}[h]
\centering\small
\caption{E4, лучшее время по ходам, мс. Конфигурация с профилировщиком долго отстаёт.}
\label{tab:e4traj}
\begin{tabular}{lcccccc}
\toprule
Ход & 1 & 2 & 3 & 4 & 5 & 6 \\
\midrule
E4 (с профилировщиком)    & 0.2999 & отказ  & 0.3050 & 0.1515 & 0.1269 & \textbf{0.1133} \\
Разбор ошибок (E1, опора) & 0.3711 & 0.1267 & 0.1261 & 0.1179 & 0.1164 & 0.1160 \\
\bottomrule
\end{tabular}
\end{table}

Информативнее итога --- траектория. На первом ходу конфигурация с профилировщиком даже
впереди, но затем проваливает второй ход и до пятого включительно остаётся позади; обгон
происходит ровно на последнем, шестом ходу. Профиль --- отложенная инвестиция: сначала
нужно получить хоть какое-то работающее ядро, потом осмыслить его разбор, и только после
этого начинается целенаправленная оптимизация вместо угадывания. Опорная конфигурация тем
временем выходит на плато рано и почти не улучшается с третьего хода (0.1261 $\to$ 0.1160).

Цена --- 11.17\,млн входных токенов против 7.43\,млн (+50\,\%). На горизонте в шесть ходов
это плохая сделка: два процента ускорения за пятьдесят процентов бюджета. Но форма кривых
говорит, что мы видим только начало --- опорная конфигурация свой потенциал исчерпала, а
конфигурация с профилировщиком на последнем ходу всё ещё улучшается на 11\,\% за ход.
Проверять это нужно увеличением числа ходов, а не повторением текущей конфигурации.

\end{document}